\begin{tikzpicture}[
    font=\small,
    cluster/.style={
        draw=black!55,
        fill=black!2,
        rounded corners=8pt,
        thick,
        inner xsep=18pt,
        inner ysep=20pt
    },
    nodebox/.style={
        draw=black!65,
        fill=black!4,
        rounded corners=7pt,
        thick,
        inner xsep=16pt,
        inner ysep=18pt
    },
    pod/.style={
        draw=blue!65!black,
        fill=blue!10,
        rounded corners=4pt,
        minimum width=2.15cm,
        minimum height=1.05cm,
        align=center,
        thick
    },
    clientpod/.style={
        draw=orange!75!black,
        fill=orange!13,
        rounded corners=4pt,
        minimum width=2.25cm,
        minimum height=1.05cm,
        align=center,
        thick
    },
    svc/.style={
        draw=green!45!black,
        fill=green!12,
        rounded corners=3pt,
        minimum width=1.75cm,
        minimum height=0.55cm,
        align=center,
        thick,
        font=\scriptsize
    },
    arrow/.style={
        -{Latex[length=2mm]},
        thick,
        draw=black!70
    },
    returnarrow/.style={
        -{Latex[length=2mm]},
        thick,
        draw=black!55,
        dashed
    },
    smalllabel/.style={
        font=\scriptsize,
        text=black!65,
        align=center
    },
    segmentlabel/.style={
        font=\scriptsize,
        text=blue!65!black,
        align=center
    },
    transferlabel/.style={
        font=\scriptsize,
        text=red!70!black,
        align=center
    },
    node distance=0.55cm
]

% Top row: client and inference pods
\node[clientpod] (client) {Benchmark\\client pod};

\node[pod, right=1.05cm of client] (p1) {Pod 1\\\texttt{stem+layer1}};
\node[pod, right=0.95cm of p1] (p2) {Pod 2\\\texttt{layer2}};
\node[pod, right=0.95cm of p2] (p3) {Pod 3\\\texttt{layer3}};
\node[pod, right=0.95cm of p3] (p4) {Pod 4\\\texttt{layer4}};
\node[pod, right=0.95cm of p4] (p5) {Pod 5\\\texttt{avgpool+fc}};

% Services below each inference pod
\node[svc, below=0.45cm of p1] (s1) {Service 1};
\node[svc, below=0.45cm of p2] (s2) {Service 2};
\node[svc, below=0.45cm of p3] (s3) {Service 3};
\node[svc, below=0.45cm of p4] (s4) {Service 4};
\node[svc, below=0.45cm of p5] (s5) {Service 5};

% Background boxes: one Kubernetes node inside one cluster
\begin{pgfonlayer}{background}
\node[
    nodebox,
    fit=(client)(p1)(p2)(p3)(p4)(p5)(s1)(s2)(s3)(s4)(s5),
    label={[font=\small\bfseries, text=black!70]above:Single Kubernetes node}
] (k8snode) {};

\node[
    cluster,
    fit=(k8snode),
    label={[font=\small\bfseries, text=black!70]above:Local Kubernetes cluster / namespace \texttt{rq14}}
] (clusterbox) {};
\end{pgfonlayer}

% Request and forwarding path
\draw[arrow] (client) -- node[smalllabel, above] {gRPC} (p1);
\draw[arrow] (p1) -- node[transferlabel, above] {784 KiB} (p2);
\draw[arrow] (p2) -- node[transferlabel, above] {392 KiB} (p3);
\draw[arrow] (p3) -- node[transferlabel, above] {196 KiB} (p4);
\draw[arrow] (p4) -- node[transferlabel, above] {98 KiB} (p5);

% Pod-to-service visual association
\draw[black!35, thick] (s1) -- (p1);
\draw[black!35, thick] (s2) -- (p2);
\draw[black!35, thick] (s3) -- (p3);
\draw[black!35, thick] (s4) -- (p4);
\draw[black!35, thick] (s5) -- (p5);

% Return path
\draw[returnarrow]
    (p5.south) .. controls +(0,-1.25) and +(0,-1.25) ..
    node[smalllabel, below] {final output returned to caller}
    (client.south);

% Explanatory labels
\node[
    smalllabel,
    below=0.55cm of s3,
    text width=9.5cm
] {
    \texttt{chain\_5svc} is the maximum coarse-stage decomposition:
    each residual-stage boundary becomes a synchronous gRPC hop between pods.
};

\node[
    segmentlabel,
    above=0.35cm of p3,
    text width=7.5cm
] {
    Forward path: one request enters the first pod and intermediate activations are forwarded synchronously until inference completes.
};

\end{tikzpicture}%